\documentclass[11pt]{article}
% Préambule commun des prétirages v2 (format arXiv : source TeX autonome, pdfLaTeX).
\usepackage[T1]{fontenc}
\usepackage[utf8]{inputenc}
\usepackage[french]{babel}
\usepackage{lmodern,microtype}
\usepackage[a4paper,margin=2.5cm]{geometry}
\usepackage{amsmath,amssymb}
\usepackage{graphicx,booktabs,siunitx}
\usepackage{authblk}
\usepackage[numbers,sort&compress]{natbib}
\usepackage[hidelinks]{hyperref}
\usepackage{csquotes}
\sisetup{locale=FR,group-minimum-digits=4}
\renewcommand\Affilfont{\small}
% Étiquette de provenance obligatoire sur chaque nombre (règle du projet) :
\newcommand{\mesure}{\textsuperscript{\textsc{m}}}   % mesuré
\newcommand{\derive}{\textsuperscript{\textsc{d}}}   % dérivé
\newcommand{\suppose}{\textsuperscript{\textsc{s}}}  % supposé


\title{TITRE : une phrase qui porte le résultat principal}
\author[1]{Équipe hackathon-santé}
\affil[1]{Projet « jumeau numérique de l'offre de soins », dépôt \texttt{hackathon-sante/v2}}
\date{Version du \today}

\begin{document}
\maketitle

\begin{abstract}
Contexte, question, méthode, résultat chiffré (avec sa provenance), limite principale. 150 à 250 mots.
\end{abstract}
\noindent\textbf{Mots-clés :} démographie des professions de santé ; projection ; \ldots

\section{Introduction}      % question, pourquoi elle compte pour la décision publique, contributions (liste)
\section{Travaux existants} % sources citées et vérifiées uniquement (DREES, Insee, IRDES, littérature)
\section{Données}           % chaque jeu : producteur, millésime, grain, licence, contrôles
\section{Méthode}           % équations, hypothèses nommées, ce que chaque levier modifie
\section{Résultats}         % chiffres du rapport v2/docs/lab/reports/, étiquetés \mesure / \derive / \suppose
\section{Validation et menaces à la validité} % protocole pré-enregistré (axe 3), ce qui pourrait rendre le résultat faux
\section{Reproductibilité}  % commandes exactes, graines, versions des données et du code
\section{Conclusion}        % ce que ça change pour les autres axes, questions ouvertes avec leur coût

\bibliographystyle{plainnat}
\bibliography{refs}
\appendix
\section{Annexes}
\end{document}
